% ===================================================================
%  TAULA N.º 2 — El Cos Humà. — L'Esbarjo. — Els Jocs.
%  Traduction 2026 d'apres texto/io, controlee sur texto/fr et
%  texto/es. Consigne : voir texto/ca/10-taula-01.tex.
% ===================================================================

%%K t02-apar-1 apar
\VUsaut{4.0mm}
\VUtitre{15.0pt}{60}{TAULA N.º 2}
\VUsaut{2.0mm}
\VUtitre{11.4pt}{0}{El Cos Humà. --- L'Esbarjo.}
\VUtitre{11.4pt}{0}{Els Jocs.}

%%K t02-tit-0 sub
\VUtitre{13.2pt}{0}{El Cos Humà.}
\VUsaut{2.4mm}
\VUcentre{10.2pt}{0}{\textit{(Lliçó senzilla de ciències naturals.)}}
\VUsaut{3.0mm}

%%K t02-tit-1 sub pos=t02-01-1
\VUpk{10.2pt}{40}{El Cap.}
\VUsaut{3.0mm}

%%K t02-01-1 p
1. --- Les principals parts del cos humà són: el \VUgras{cap}, el \VUgras{tronc}
i els \VUgras{membres}.

%%K t02-02-1 p
2. --- El cap consta de dues parts: el \VUgras{crani}, que conté el
\VUgras{cervell} i que cobreixen els \VUgras{cabells}, i la \VUgras{cara}.

%%K t02-03-1 p
3. --- En el nostre dibuix no veiem ni el crani ni el cervell; però veiem amb
tota claredat les parts importants de la cara.

%%K t02-04-1 p
4. --- Heus aquí primer el \VUgras{front}, que anuncia una intel·ligència
poderosa, un pensament sempre despert. Els \VUgras{polsos} són molt lliures.
L'\VUgras{ull} és viu i ben obert. Una \VUgras{cella} espessa i unes
\VUgras{pestanyes} llargues el protegeixen.

%%K t02-05-1 p
5. --- Heus aquí a més el \VUgras{nas} i els \VUgras{narius}. El nas ens serveix
per olorar i per respirar. A cada costat del nas hi ha les \VUgras{galtes}, i més
enllà les \VUgras{orelles}. La part inferior de la cara la formen la
\VUgras{boca} i la \VUgras{barbeta}.

%%K t02-06-1 p
6. --- No les veiem gaire bé, perquè el \VUgras{bigoti} i les \VUgras{patilles}
ens les amaguen. Però sabem que la boca té els dos \VUgras{llavis}, la
\VUgras{mandíbula superior} i la \VUgras{mandíbula inferior} amb les seves
\VUgras{dents}, i la \VUgras{llengua}. Amb la boca parlem i mengem. Amb la
llengua assaborim els aliments.

%%K t02-07-1 p
7. --- L'ull, l'orella, el nas i la boca són, juntament amb els dits, els òrgans
dels cinc sentits: la \VUgras{vista}, l'\VUgras{oïda}, l'\VUgras{olfacte}, el
\VUgras{gust} i el \VUgras{tacte}.

%%K t02-08-1 p
8. --- El cap és, doncs, una de les parts més interessants del cos humà.

%%K t02-tit-2 sub
\VUsaut{4.4mm}
\VUpk{10.2pt}{40}{El Tronc.}
\VUsaut{3.0mm}

%%K t02-09-1 p
9. --- El \VUgras{coll} uneix el cap al tronc. Comprèn la \VUgras{gola} i el
\VUgras{clatell}.

%%K t02-10-1 p
10. --- El tronc conté també molts òrgans essencials. Al \VUgras{pit} hi ha els
\VUgras{pulmons}, amb els quals respirem, i el \VUgras{cor}, que és el centre de
la vida. Quan el cor deixa de bategar, l'ésser viu mor.

%%K t02-11-1 p
11. --- Les \VUgras{costelles} sostenen el pit.

%%K t02-12-1 p
12. --- Les altres parts del tronc són el \VUgras{costat}, l'\VUgras{esquena},
les \VUgras{espatlles} i el \VUgras{ventre}. Ens ajaiem de costat o d'esquena per
dormir, i portem les càrregues damunt l'espatlla.

%%K t02-13-1 p
13. --- Al ventre hi ha l'\VUgras{estómac} i els \VUgras{intestins}. Ens
serveixen per digerir els aliments.

%%K t02-tit-3 sub
\VUsaut{4.4mm}
\VUpk{10.2pt}{40}{Els Membres.}
\VUsaut{3.0mm}

%%K t02-14-1 p
14. --- Tenim quatre \VUgras{membres}: dos \VUgras{braços} i dues
\VUgras{cames}. Amb els braços agafem, sostenim, aixequem i recollim els
objectes; amb les cames ens tenim drets, caminem i correm.

%%K t02-15-1 p
15. --- L'\VUgras{espatlla} uneix el braç al cos. Un \VUgras{múscul} molt
robust, el bíceps, mou el braç, que el \VUgras{colze} uneix a l'\VUgras{avantbraç}.
El \VUgras{canell} forma l'articulació de la \VUgras{mà}, que acaba en els
\VUgras{dits} i les \VUgras{ungles}. A la mà distingim una \VUgras{vena} (i una
artèria) per la qual corre la \VUgras{sang}.

%%K t02-16-1 p
16. --- Les parts de la cama són: la \VUgras{cuixa}, el \VUgras{genoll} i la
\VUgras{sofraja}, el \VUgras{panxell}, la \VUgras{carn} del qual cobreix la
\VUgras{pell}, el \VUgras{turmell} i el \VUgras{peu}. Els \VUgras{ossos} de la
cama són molt gruixuts i molt forts. A l'extrem del peu hi ha els \VUgras{dits
del peu}. Les altres parts són la \VUgras{planta} i el \VUgras{taló}.

%%K t02-17-1 p
17. --- Tal és la descripció gairebé completa del cos humà.

%%K t02-17-2 p
\textit{Nota.} --- \textit{Amb l'ajuda de l'expressió «em fa mal... (el cap,
etc.)», el professor podrà repassar tot aquest vocabulari des del punt de vista
del bon o del mal} \VUgras{estat del cos}.

%%K t02-tit-4 sub
\VUsaut{5.0mm}
\VUpk{10.2pt}{40}{Gimnàstica.}
\VUsaut{2.4mm}
\VUcentre{10.2pt}{0}{\textit{(Contat per un dels alumnes.)}}
\VUsaut{3.0mm}

%%K t02-18-1 p
18. --- Per ser flexibles, àgils i sans, hem d'exercitar les cames i els braços.
Per això juguem i fem \VUgras{gimnàstica}.

%%K t02-19-1 p
19. --- Entre setmana, aquests dos exercicis es fan al pati de l'institut (vegeu
la taula n.º 1). Però els dijous i els diumenges anem a una gran \VUgras{prada},
on tenim lloc per córrer i divertir-nos.

%%K t02-20-1 p
20. --- Els nostres mestres i els nostres pares ens hi acompanyen de vegades;
però és el \VUgras{professor de gimnàstica} \textsuperscript{(12)} qui dirigeix
els exercicis.

%%K t02-21-1 p
21. --- A la prada, a l'esquerra de l'entrada, hi ha els \VUgras{aparells de
gimnàstica} \textsuperscript{(1)}: enfilem-nos per l'\VUgras{escala}
\textsuperscript{(2)}, fem el pi a les \VUgras{anelles} \textsuperscript{(3)} i
al \VUgras{trapezi} \textsuperscript{(4)}, ens hissem amb les mans fins dalt de
l'\VUgras{escala de corda} \textsuperscript{(5)} o de la \VUgras{corda de nusos}
\textsuperscript{(6)}.

%%K t02-22-1 p
22. --- Mentrestant, els nostres companys salten per damunt del \VUgras{cavall
de fusta} \textsuperscript{(7)} o de les \VUgras{barres paral·leles}
\textsuperscript{(11)}. D'altres \VUgras{boxegen} \textsuperscript{(8)} o fan
\VUgras{esgrima} \textsuperscript{(9)} amb careta i \VUgras{floret}. D'altres,
finalment, aixequen \VUgras{peses} \textsuperscript{(10)}.

%%K t02-tit-5 sub
\VUsaut{4.6mm}
\VUpk{10.2pt}{40}{Els Jocs.}
\VUsaut{3.0mm}

%%K t02-23-1 p
23. --- Al voltant dels gimnastes, nois i noies juguen a diferents
\VUgras{jocs}. Les noies es diverteixen amb el seu \VUgras{cèrcol}
\textsuperscript{(14)}; salten a la \VUgras{corda} \textsuperscript{(13)}, o
conviden els seus germans i els seus cosins a una partida de \VUgras{croquet}
\textsuperscript{(15)}. Els encanta apartar la \VUgras{bola} del contrari amb la
seva \VUgras{maça de fusta}, just quan aquest es pensava passar sense treball per
l'arquet!

%%K t02-24-1 p
24. --- Els nois prefereixen els jocs més recis. Juguen de bona gana a les
\VUgras{bitlles} \textsuperscript{(16)}, que abaten amb traça amb la
\VUgras{bola} \textsuperscript{(17)}. Tampoc no menyspreen la \VUgras{baldufa}
\textsuperscript{(18)} i el \VUgras{bilboquet} \textsuperscript{(19)}, ni tan
sols el \VUgras{joc de la granota} \textsuperscript{(20)}. Però els seus jocs
preferits són les \VUgras{bales} \textsuperscript{(21)} i la \VUgras{pilota a la
paret} \textsuperscript{(22)}, perquè la paret de l'\VUgras{hivernacle}
\textsuperscript{(26)} va molt bé per fer rebotar la pilota lleugera.

%%K t02-25-1 p
25. --- El petit Joan, enfilat a les seves \VUgras{xanques} \textsuperscript{(24)},
és molt destre i sap guardar molt bé l'equilibri. Prop d'ell, alguns companys
juguen a \VUgras{fet i amagar} \textsuperscript{(25)} en aquell hivernacle i
darrere la \VUgras{tanca}. D'altres prefereixen la \VUgras{bufetada endevinada}
\textsuperscript{(28)} o el \VUgras{volant} \textsuperscript{(29)}, que voleia
d'una \VUgras{raqueta} a l'altra.

%%K t02-noto-1 noto
\VUnotes{143.2mm}{%
(*) Els jocs infantils tenen sovint noms purament nacionals. Els hem anomenat en
ido tan bé com hem pogut, ja inspirant-nos en l'acció mateixa que els
caracteritza, ja adoptant el nom nacional d'un o altre país quan no resulta massa
inexacte. Per exemple: el \VUgras{joc de la bóta} (alemany i francès) és el
\VUgras{joc de la granota} \textsuperscript{(20)} (anglès), en què els jugadors
proven de tirar els seus discos rodons per la boca de la granota de metall que
és amb la boca oberta damunt la caixa (la bóta) foradada. El \VUgras{salt a
l'espatlla} \textsuperscript{(30)} es distingeix del \VUgras{salt a l'esquena}
només en això: per saltar per damunt del cap, els jugadors recolzen les mans a
les espatlles del company, mentre que en el «\VUgras{salt a l'esquena}» les
recolzen només a la seva esquena. --- O bé, alguns dels qui juguen s'inclinen
mentre els altres salten damunt les seves esquenes recolzant les mans a
l'espatlla; els primers han d'aguantar el cop sense doblegar-se, els segons han
de saltar sense perdre l'equilibri (vegeu la taula mateixa).%
}

%%K t02-26-1 p
26. --- Al mig de la prada hi ha dos arbres. Al peu d'un d'ells, set o vuit nois
han organitzat una partida de \VUgras{salt a l'espatlla} \textsuperscript{(30)}
(*). Aquest joc no es coneix a tots els països; però tothom sap què és el
\VUgras{salt a l'esquena}, al qual els nois no juguen ara.

%%K t02-tit-6 sub
\VUsaut{5.0mm}
\VUpk{10.2pt}{40}{Els Jocs a l'aire lliure.}
\VUsaut{3.4mm}

%%K t02-27-1 p
27. La \VUgras{gallina cega} \textsuperscript{(31)} és també molt divertida;
noies i nois es posen per torn la \VUgras{bena} als ulls i atrapen els maldestres
que es deixen agafar.

%%K t02-28-1 p
28. --- Al mig de la prada, heus aquí un joc molt francès, encara que una mica
brusc: és l'\VUgras{óssa} \textsuperscript{(32)}, molt més sorollosa que el joc
tan tranquil de l'\VUgras{estel} \textsuperscript{(33)} o que el joc anglès del
\VUgras{tennis} \textsuperscript{(34)}.

%%K t02-29-1 p
29. --- Aquest darrer esport fa les delícies dels alumnes grans. Els nostres
propis professors el practiquen de bona gana. Exigeix molta destresa i agilitat,
i a mi m'agrada moltíssim. Deixo el \VUgras{gronxador} \textsuperscript{(35)}
per a les noies.

%%K t02-30-1 p
30. --- Tampoc no m'agrada un altre joc anglès que pot arribar a ser perillós.

%%K t02-30-1b p
Em refereixo al \VUgras{futbol} \textsuperscript{(36)}, que fa la felicitat del
meu germà gran. Prefereixo el nostre vell \VUgras{joc de barres}
\textsuperscript{(38)}, igual de sa i molt menys brutal.

%%K t02-tit-7 sub
\VUsaut{4.6mm}
\VUpk{10.2pt}{40}{Ciclisme i jocs de pilota.}
\VUsaut{3.0mm}

%%K t02-31-1 p
31. --- També dono amb gust una volta a la prada en \VUgras{bicicleta}
\textsuperscript{(37)}.

%%K t02-32-1 p
32. --- L'any que ve aprendré a \VUgras{muntar a cavall}
\textsuperscript{(39)}. Que n'és, de bonic, un cavall que camina o trota amb
gràcia, i que salta els obstacles al galop!

%%K t02-33-1 p
33. --- A l'estiu aprenem a \VUgras{nedar} \textsuperscript{(40)}. Ens capbussem
i ens banyem amb delit a l'aigua clara i fresca del riu. De vegades, al final,
deixem anar un \VUgras{globus de paper} \textsuperscript{(41)} que puja
majestuosament per l'aire tan aviat com s'omple d'aire calent.

%%K t02-34-1 p
34. --- Hi ha encara molts altres jocs, però no acabaria mai d'anomenar-los, i
per aquest breu resum es veu que els dijous a la nostra bella prada no són gens
avorrits.

%%K t02-34-2 p
N. B. --- \textit{El professor completarà fàcilment aquest capítol descrivint
amb més detall cadascun dels jocs importants.}
\VUsaut{6.0mm}
\VUfilet{22mm}
